%% TODO 
%% Elman's temporal XOR example: https://onlinelibrary.wiley.com/doi/epdf/10.1207/s15516709cog1402_1
%% Minsky's theorem  https://www.dlsi.ua.es/~mlf/nnafmc/pbook/node16.html
%% Siegelmann and Sontag results: https://reader.elsevier.com/reader/sd/pii/089396599190080F?token=3BB6C8998A5747F205F88497E16A35E12668BE9E0961223F1B6DF756D111BD3833E569B62DB794C103A71A11C2496066&originRegion=eu-west-1&originCreation=20221115155259
%% https://www.sciencedirect.com/science/article/pii/S0022000085710136


\subsection{Recurrent Neural Networks} \label{sec:rnns}
Natural languages are beyond the descriptive power of regular language. Thus, finite-state models are not sufficient for natural language modeling.
Now, we turn to a class of models that is more powerful than finite-state models. 

Recurrent neural networks (RNNs) have been used to process sequential data and time series. 
In this subsection, we focus on their application to modeling languages.

The most unique characteristic of RNNs is that they take previous outputs as input. Such recurrent formulation endows them with the property of ``unbounded dependency'' (thereby ``unbounded memory'') on the input data, distinguishing them from the Markovian \ngram{} models. 


This recurrent formulation requires sharing parameters through time --- by repeatedly applying the same update rule to the hidden state of the model.
The update rule of RNNs is called \defn{dynamics map}:

\begin{definition}
A \defn{dynamics map} $f$ is a function of the type $\Rd \times \alphabet \rightarrow \Rd$. 
\end{definition}

\begin{definition}\label{def:recurrent-neural-language-model}
A $d$-dimensional \defn{recurrent neural language model} over alphabet $\alphabet$ with dynamics map $f$ is a distribution of the form
\begin{equation}
p(\vy) \defeq \prod_{t=1}^T p(y_t \mid \vy_{<t}) = \prod_{t=1}^T \frac{\exp\left(\embedSymt^{\top} \vht + \biasSymt \right) }{\sum_{y' \in \eosalphabet} \exp\left(\embedSymPrime^{\top} \vht+ \biasSymPrime\right)} 
\end{equation}
where $\eosalphabet \defeq \alphabet \cup \{\eos\}$ and $\embedSymt, \vht \in \Rd$ and $\biasSymt \in \R$.
The hidden state is recurrently updated according to the dynamics map as follows:
\begin{align}
    \vhzero &= \mathbf{0} \\
    \vht    &= f(\vhtminus, y_{t-1})
\end{align}
\end{definition}




\subsection{General Results on Tightness}
We now prove general results on the tightness of recurrent neural language models, as defined in \Cref{def:recurrent-neural-language-model}.
The analysis is straightforward and is a translation of the generic results on tightness to the case of the norm of $\vht$.
\begin{theorem}
A recurrent neural language model is tight if and only if $||\vht||_2 \leq \cdots$.
\end{theorem}


\subsection{Elman and Jordan Networks}
Next, we introduce two of the simplest recurrent neural language models.
We term them Elman language models and Jordan language models, as each is inspired by architectures proposed by \citet{Elman1990} and \citet{Jordan1986}, respectively. 
The definitions we offer are slightly different than those found in the original works---most notably, both Elman and Jordan construct neural networks for transduction, rather than language modeling

\begin{definition}[Elman Network]
An \defn{Elman language model} with non-linearity $\sigma$ is a $d$-valued recurrent neural language model over an alphabet $\alphabet$ with the following dynamics map\looseness=-1
\begin{align}
\vht &= \sigmoid\left(\rmU \mathhl{\vhtminus} + \rmW \embedSymtminus + \biasVech \right) \\
\vot &= \sigmoid\left(\rmV \vhtminus + \biasVecy \right)
\end{align}
where the $\sigmoid$ is applied element-wise and $\rmU,\rmV,\rmW \in \R^{d \times d}$. 
\end{definition}

\begin{definition}[Jordan Network]
A \defn{Jordan language model} with non-linearity $\sigma$ is a $d$-valued recurrent neural language model over an alphabet $\alphabet$ with the following dynamics map\looseness=-1
\begin{align}
\vht &= \sigmoid\left(\rmU \mathhl{\votminus} + \rmW \embedSymtminus + \biasVech \right) \\
\vot &= \sigmoid\left(\rmV \vhtminus + \biasVecy \right)
\end{align}
where the $\sigmoid$ is applied element-wise and $\rmU,\rmV,\rmW \in \R^{d \times d}$. 
\end{definition}

\begin{definition}
The \defn{Heaviside} function is defined as
\begin{equation}\label{eq:heaviside}
H(x) = \begin{cases} 1 & \ifcondition x > 0 \\
0 & \otherwisecondition
\end{cases}
\end{equation}
\end{definition}


\begin{lemma}
The distribution represented by a $d$-dimensional Jordan network with a Heaviside non-linearity $H$ is regular.
\end{lemma}
\begin{proof}
We construct a deterministic probabilistic FSA that represents the same distribution.
First, we note that, by the definition of the Heaviside function (\Cref{eq:heaviside}) every $\vht \in \{0, 1\}^d$. 
We construct a set of $2^d$ states $\states$, each of which is associated with a binary vector. 
Now, for every state $q \in \states$, add a transition $q \xrightarrow[]{a / w} q'$ where $q' = f(q, w)$ and 
\ryan{add definition}
\end{proof}

\begin{example}
\ryan{Tianyu, can you write down a simple example?}
\end{example}

\begin{example}
\ryan{Tianyu, can you write down a more complicated example? I want this example to express a distribution that would require an exponential number of states for a PFSA.}
\end{example}

\begin{lemma}
Any regular language can be encoded as an Jordan network.
\end{lemma}


\begin{theorem}
Simple recurrent neural language models are as expressed as probabilistic FSAs.\ryan{I think we might need deterministic for this to be true}
\end{theorem}
 

\paragraph{A Motivating Example.}
The need for memory in neural networks arises especially when we are dealing with temporal problems. The recurrence enables the neural networks to model patterns that repeats overtime. 

\begin{example}[Relative temporal relations]
    
    
\end{example}
\ryan{Tianyu, can you add the XOR example here.}

\paragraph{Tightness of Elman Networks.}
\Ryan{Example: Exhibit a non-tight language model}

\Ryan{Theorem: Any Elman network with a bounded activation function is tight}


\subsubsection{Analysis: The Vanishing/Exploding Gradient Problem}
The vanishing and exploding gradient are common problems in training RNNs with backpropagation. 

Taking the Elman language model as an example, 
recall the hidden state update of recurrent neural networks:
\begin{align}
    \vhzero &= \mathbf{0} \\
    \vht    &= \func(\vhtminus, y_{t-1}) = \sigmoid \left( \rmU \vhtminus + \rmW \embedSymtminus + \biasVech \right)
\end{align}

The step-wise gradient for one step of hidden state update is:
\begin{align}
    \frac{\partial \vht}{\partial \vhtminus} &= \frac{\partial \func(\vhtminus, y_{t-1})}{\partial \vhtminus} \\
    &= \diag{ \sigmoid^\prime \left( \rmU \vhtminus + \rmW \embedSymtminus + \biasVech \right)} \rmU
\end{align}

The problems appear when the input sequence grows longer. 
By applying the chain rule, we obtain the effect of $\vhone$ on $\vhT$: 
\begin{align}
    \frac{\partial \vhT}{\partial \vhone} 
    &= \prod_{t=1}^{T} \frac{\partial \vht}{\partial \vhtminus} \\
    &= \prod_{t=1}^{T} \diag{\sigmoid^\prime \left( \rmU \vhtminus + \rmW \embedSymtminus + \biasVech \right)} \rmU
\end{align}

We denote $\norm{\rmU}$ as the $\normltwo$ norm of $\rmU$. 
The vanishing gradient problem appears when $\norm{\rmU} < 1$:
\begin{align}
   \norm{\frac{\partial \vhT}{\partial \vhone}}
    &\le \prod_{t=1}^{T} \norm{\frac{\partial \vht}{\partial \vhtminus}} \\
    &\le \norm{\rmU}^T \prod_{t=1}^{T} \norm{\diag{\sigmoid^\prime \left( \rmU \vhtminus + \rmW \embedSymtminus + \biasVech \right)}} 
\end{align}

$\norm{\rmU} < 1$ is the \emph{sufficient} condition for vanishing gradient. The term $\norm{\rmU}^T$ exponentially shrinks to $0$ when $\norm{\rmU} < 1$, which results in a vanishing gradient. 


The proof is a bit different for exploding gradient. A weaker result can be derived by inverting the proof for vanishing gradient:

$\norm{\rmU} \geq 1$ is the \emph{necessary} condition for exploding gradient. Since otherwise, the gradient would vanish. 


\paragraph{Why is vanishing gradient a problem?} 
Gradient signals are lost for faraway inputs, since they are far smaller in scale than the signals for nearby inputs. 
Gradient signals could also be rounded off to zero due to the finite float precision. 

\paragraph{Why is exploding gradient a problem?} 
A gradient that is too large in scale could lead to $\textnormal{NaN}$ updates when training an RNN. 


The vanishing/exploding gradient problem can be alleviated by some tricks during training. 
\begin{itemize}
    \item Gradient clipping. Clipping gradients to some customized range. 
    \item Gradient normalization. Normalizing gradients by their norm.
    \item \ldots
\end{itemize}


It also motivates some RNN variants with more stable gradients on long sequence inputs, such as LSTM and GRU. 


\subsubsection{Analysis: Thresholding Activations}
One fruitful manner to analyze the expressivity of recurrent neural networks is by making simplifying assumptions on the non-linear activations.
A common simplification that goes back to the 1960s is the thresholding sigmoid, which we define below.
\begin{definition}[Thresholding Sigmoid]
$$
    \sigma_{\textnormal{T}}(x)=
    \begin{cases}
    0 &  \ifcondition x  \le 0 \\
    1 &  \ifcondition x  > 0
    \end{cases}
$$
\end{definition}
In words, the thresholding function maps every real value to a $0$ or to a $1$, depending on whether it is greater than or less than zero. 

Here we give the definition of the thresholding Elman network, 
\begin{definition}[Thresholding Elman Network]
A \defn{Thresholding Elman network} is an Elman network with thresholding sigmoid $ \sigma_{\textnormal{T}}$ as the non-linearity. 
\end{definition}

We now give a variant of a classic theorem originally due to \citet{Minsky1954} but with a probabilistic twist.


\begin{theorem}
Let $\wfsa = \wfsatuple$\ryan{Insert definition} be a tight probabilistic deterministic finite-state automaton.
Then, there exists a thresholding Elman network\ryan{This should be introduced} with a hidden state of size $|\alphabet||\states|$ that encodes the same distribution as $\wfsa$.
\end{theorem}
\begin{proof}
We give proof by construction.
Each hidden state lives in $\vh \in \mathbb{B}^{|\alphabet||\states|}$. \ryan{macro the boolean set} due to the thresholding function.
Let $f : \states \times \alphabet \rightarrow [|\states||\alphabet|]$ be a bijection.
\Ryan{Mention the semantics}

The matrix\ryan{Maybe give it a name} $U$ lives in $\R^{|\alphabet||\states| \times |\alphabet||\states|}$.
Let $U_{ik} = 1$ if $i = f(q, b)$ and $k = f(p, a)$ and there exists a state $p$ such that $p\xrightarrow[b]{q}$ is a transition in $\automaton$; otherwise, let $U_{ik} = -1$. 
Because $a$ is free, we have that each row has $|\alphabet|$ entries that are $1$; the rest are $-1$. 
The matrix $V$ lives in $\R^{|\alphabet||\states| \times |\alphabet|}$.
Let $V_{ik} = 1$ 



\end{proof}

\subsection{Analysis: Saturated Activations}



\begin{definition}[Saturated Sigmoid]
$$
\sigma(x)=
\begin{cases}
0 &  \ifcondition x  \le 0  \\
x &  \ifcondition 0 < x \le  1 \\
1 &  \ifcondition x  > 1 
\end{cases}
$$
\end{definition}


\begin{lemma}
Consider an 
\end{lemma}



\subsubsection{Long Short-term Memory}

\subsubsection{Gated Recurrent Units}